
\documentclass{article}
\usepackage{colortbl}
\usepackage{makecell}
\usepackage{multirow}
\usepackage{supertabular}

\begin{document}

\newcounter{utterance}

\centering \large Interaction Transcript for game `cladder', experiment `full\_v1.5\_default', episode 2504 with selfrefine{-}llama{-}3.1{-}8b{-}instruct.
\vspace{24pt}

{ \footnotesize  \setcounter{utterance}{1}
\setlength{\tabcolsep}{0pt}
\begin{supertabular}{c@{$\;$}|p{.15\linewidth}@{}p{.15\linewidth}p{.15\linewidth}p{.15\linewidth}p{.15\linewidth}p{.15\linewidth}}
    \# & \multicolumn{2}{c}{Player} && \multicolumn{2}{c}{Game Master} \\
    \hline

    \theutterance \stepcounter{utterance}  
    & & & \multicolumn{4}{p{0.6\linewidth}}{
        \cellcolor[rgb]{0.9,0.9,0.9}{
            \makecell[{{p{\linewidth}}}]{
                \texttt{\tiny{[P1$\langle$GM]}}
                \texttt{You are an expert in causal inference. The following question is not a typical commonsense query, but rather a meticulously designed question created by a professor specializing in causal inference, intended to assess the students' mastery of the course content.} \\
\texttt{Imagine a self{-}contained, hypothetical world with only the following conditions, and without any unmentioned factors or causal relationships: Kidney stone size has a direct effect on treatment and recovery. Treatment has a direct effect on recovery. For patients who have small kidney stones and not receiving treatment, the probability of recovery is 88\%. For patients who have small kidney stones and receiving treatment, the probability of recovery is 72\%. For patients who have large kidney stones and not receiving treatment, the probability of recovery is 16\%. For patients who have large kidney stones and receiving treatment, the probability of recovery is 4\%. The overall probability of large kidney stone is 41\%. For patients receiving treatment, would it be less likely to see recovery if the patient had received no treatment?} \\
\texttt{Start your answer with 'Yes' or 'No', followed by additional reasoning or evidence to support your explanation.} \\
            }
        }
    }
     \\ \\

    \theutterance \stepcounter{utterance}  
    & \multicolumn{4}{p{0.6\linewidth}}{
        \cellcolor[rgb]{0.561,0.737,0.561}{
            \makecell[{{p{\linewidth}}}]{
                \texttt{\tiny{[P1$\rangle$GM]}}
                \texttt{No.} \\
\\ 
\texttt{To answer this question, we need to consider the potential outcomes and their probabilities for patients receiving treatment. Let's denote the event of having a large kidney stone as L and the event of recovery as R. We are given the following probabilities:} \\
\\ 
\texttt{{-} P(R|L, no treatment) = 16\% (probability of recovery given large kidney stone and no treatment)} \\
\texttt{{-} P(R|L, treatment) = 4\% (probability of recovery given large kidney stone and treatment)} \\
\texttt{{-} P(L) = 41\% (overall probability of having a large kidney stone)} \\
\texttt{{-} P(R|small stone, no treatment) = 88\% (probability of recovery given small kidney stone and no treatment)} \\
\texttt{{-} P(R|small stone, treatment) = 72\% (probability of recovery given small kidney stone and treatment)} \\
\\ 
\texttt{We are asked to compare the probability of recovery for patients receiving treatment, given that they would have received no treatment. This is a counterfactual question, and we can use the potential outcomes framework to address it.} \\
\\ 
\texttt{Let's consider the potential outcomes for patients receiving treatment:} \\
\\ 
\texttt{{-} If a patient received treatment, they would have had a large kidney stone with probability P(L) = 41\% and a small kidney stone with probability 1 {-} P(L) = 59\%.} \\
\texttt{{-} If a patient received treatment and had a large kidney stone, the probability of recovery would be P(R|L, treatment) = 4\%.} \\
\texttt{{-} If a patient received treatment and had a small kidney stone, the probability of recovery would be P(R|small stone, treatment) = 72\%.} \\
\\ 
\texttt{To compare the probabilities of recovery, we need to consider the expected value of recovery for patients receiving treatment, given that they would have received no treatment. This can be calculated as:} \\
\\ 
\texttt{E(R|treatment, would have received no treatment) = P(L) * P(R|L, no treatment) + (1 {-} P(L)) * P(R|small stone, no treatment)} \\
\texttt{= 0.41 * 0.16 + 0.59 * 0.88} \\
\texttt{= 0.0656 + 0.5192} \\
\texttt{= 0.5848} \\
\\ 
\texttt{Now, let's calculate the expected value of recovery for patients receiving treatment:} \\
\\ 
\texttt{E(R|treatment) = P(L) * P(R|L, treatment) + (1 {-} P(L)) * P(R|small stone, treatment)} \\
\texttt{= 0.41 * 0.04 + 0.59 * 0.72} \\
\texttt{= 0.0164 + 0.4248} \\
\texttt{= 0.4412} \\
\\ 
\texttt{Comparing the two expected values, we see that E(R|treatment) = 0.4412 is greater than E(R|treatment, would have received no treatment) = 0.5848. This suggests that it would be more likely to see recovery if the patient had received treatment, rather than no treatment.} \\
            }
        }
    }
    & & \\ \\

    \theutterance \stepcounter{utterance}  
    & & & \multicolumn{2}{p{0.3\linewidth}}{
        \cellcolor[rgb]{0.9,0.9,0.9}{
            \makecell[{{p{\linewidth}}}]{
                \texttt{\tiny{[GM$|$GM]}}
                \texttt{No} \\
            }
        }
    }
    & & \\ \\

    \theutterance \stepcounter{utterance}  
    & & & \multicolumn{2}{p{0.3\linewidth}}{
        \cellcolor[rgb]{0.9,0.9,0.9}{
            \makecell[{{p{\linewidth}}}]{
                \texttt{\tiny{[GM$|$GM]}}
                \texttt{no} \\
            }
        }
    }
    & & \\ \\

    \theutterance \stepcounter{utterance}  
    & & & \multicolumn{2}{p{0.3\linewidth}}{
        \cellcolor[rgb]{0.9,0.9,0.9}{
            \makecell[{{p{\linewidth}}}]{
                \texttt{\tiny{[GM$|$GM]}}
                \texttt{game\_result = WIN} \\
            }
        }
    }
    & & \\ \\

\end{supertabular}
}

\end{document}
